\chapter{Introduction}
\label{ca:chap:intro}

\section{Motivation}
\label{ca:sec:motivation}

Much of the knowledge a practitioner would like to train on is no longer available as
data. Hospitals, banks and firms rarely release their datasets, but they increasingly
expose models trained on them, through an interface or as open weights. For anyone
training a small language model, the resource at hand is therefore less often a dataset
and more often a collection of models to query.

This project started from FedTOPS \autocite{chen_fedtops_2026}, the first study of
on-policy distillation in a federated setting. It found that a student's progress
changes what it can next learn from its teacher: averaging client updates slowed a
prerequisite skill, which delayed the point at which the teacher's supervision became
useful. FedTOPS has one teacher, reads its token probabilities, and treats every query as
free. This project asks what happens when those three assumptions are dropped. A student
can query several black-box models, each answers in text only, and each query has a price.
If the student's progress changes what it can learn from a teacher, it should also change
which teacher is worth paying for.

The obvious rule, to always query the strongest model, is known to be a poor one for
small students. Models of 3B parameters or fewer often learn more from simpler supervision
than from stronger reasoners \autocite{li_learnability_gap_2025}, and a larger teacher
does not necessarily yield a better student \autocite{guo_lada_2026}. The project
therefore studies the following question.

\begin{mdframed}[backgroundcolor=bg, linewidth=0pt, innertopmargin=8pt,
  innerbottommargin=8pt]
Given a fixed query budget and a set of black-box providers whose usefulness is unknown,
domain-dependent and changing, how should a student model decide whom to ask, about
what, and when?
\end{mdframed}

This setting is called provider-selective distillation. Providers return text only,
their training data never leaves them, none is assumed best everywhere, each query has a
price, and the value of a provider depends on the student's current state, which the
selector's own choices change. The last property separates the problem from online data
mixing \autocite{albalak_odm_2023}, multi-oracle active learning
\autocite{donmez_iethresh_2009} and LLM routing \autocite{ong_routellm_2025}: here,
learning whom to ask and learning from the answer are coupled.

\section{Problem Statement}
\label{ca:sec:problem}

A student policy $\pi_\theta$ can query $K$ providers $M_1, \dots, M_K$. Provider $M_k$
is a sampling oracle that returns a text response to a prompt at price $c_k$. Prompts come
from $D$ domains $\mathcal{X}_1, \dots, \mathcal{X}_D$. The student's utility
$U(\theta)$ is its mean accuracy over held-out sets, one per domain.

Training proceeds in rounds. In round $t$ a selector chooses an arm $a_t = (k_t, d_t)$,
a provider and a domain. A batch of $b$ prompts from $\mathcal{X}_{d_t}$ is answered by
$M_{k_t}$, the student is updated, and the budget is charged $b\,c_{k_t}$. The selector
sees the student's accuracy on a selection set, disjoint from the test sets. Its goal is
\[
  \max_{\pi_{\mathrm{sel}}} \; \mathbb{E}\big[\, U(\theta_\tau) \,\big]
  \quad \text{subject to} \quad \sum_{t=1}^{\tau} b\, c_{k_t} \le B .
\]
The quantity a selector must estimate is the marginal value of an arm per unit cost at the
current student state,
\[
  v(k, d; \theta) =
  \frac{\mathbb{E}\big[\, U\big(\operatorname{Update}(\theta, S_{k,d})\big) - U(\theta) \,\big]}
       {b\, c_k},
\]
where $S_{k,d}$ is a batch from domain $d$ answered by provider $k$. If $v$ did not
depend on $\theta$, this would be a bandit problem with costs. Because every update
changes $\theta$, rewards drift as a result of the selector's own choices.

The problem is only worth solving if provider value really varies in this way. Three
hypotheses are therefore tested in order.

\begin{description}[itemsep=2pt,topsep=4pt]
  \item[H1 (Non-dominance).] At a fixed student state, the provider of highest marginal
    value differs across domains.
  \item[H2 (State dependence).] Within a domain, the provider of highest marginal value
    changes during training.
  \item[H3 (Adaptivity pays).] An adaptive selector reaches higher test utility than the
    best static allocation at equal budget.
\end{description}

H1 alone would justify a fixed allocation per domain. H2 is what makes adaptive selection
necessary, and H3 is what makes it worthwhile. A negative result on H2 is also useful,
since it would show that an allocation fixed before training is enough.

\section{Objectives and Scope}
\label{ca:sec:objectives}

The objectives are to:
\begin{enumerate}[itemsep=1pt,topsep=3pt]
  \item build a reproducible testbed: a small open student, black-box providers that
    differ in domain, size and price, one deliberately degraded provider, and tasks with
    checkable answers;
  \item measure marginal value directly, by training from shared student checkpoints on
    each arm, to test H1 and H2;
  \item design budget-aware selectors and compare them with uniform, single-provider and
    best-static allocations to test H3;
  \item explain the results: where learning happens, what the selectors spend on, and
    why.
\end{enumerate}

Strategic providers, providers that learn from the student, formal privacy guarantees and
logit-level distillation are out of scope. Tasks are single-turn with verifiable answers,
so rewards can be computed cheaply and without a learned judge. Provider responses are
generated once and cached, and querying is simulated from the cache.
